% !PRINTED 544
\weberpage{544}
There exist 63 negative discriminants with the property that each genus contains only a single class; that there are no more, indeed even that their number is only a finite one, can of course so far only be inferred by induction, not proved. The majority of these discriminants, which have already been collected in \S\,139, are $\equiv 1, 2\ (\mathrm{mod}\ 4)$ and without square divisor, or $\equiv 8\ (\mathrm{mod}\ 16)$.

For the former the \gesperrt{class invariants} can be completely computed by formula (10), and a similar formula, into whose formation we do not wish to enter here, also leads to the goal for the discriminants divisible by 8.

For the isolated discriminants of this kind which do not fit in here, the class invariants $f(\omega)$ can be computed by one of the other methods.\footnote{Cf.\ the author's memoir: \glqq Zur komplexen Multiplikation elliptischer Funktionen\grqq{}. Mathematische Annalen, Bd.\,23. I here call attention to an error in the \gesperrt{Gauss} table which I found on the occasion of these computations: \gesperrt{Gauss'} Werke, Bd.\,II, S.\,475, the positive determinant 136 must have the designation IV, 2, not IV, 1.}

In order to carry out the calculation on a simple example that is easy to survey, we choose $m=105=3.\,5.\,7$. If we omit the values of $\delta$ which are $\equiv 1\ (\mathrm{mod}\ 8)$, since these give no contribution in the sum of the formulae (10), then we have to put:
\[
\begin{aligned}
\delta &= -3, & 5, & & 21, & & -35, \\
\delta' &= 140, & -84, & & -20, & & 12
\end{aligned}
\]
and obtain, since $g=8$:
% [sic] The source prints radicands 84 and 20 in rows 2 and 3; the positive discriminants 5 and 21 are required by (11) and are used explicitly on p.545.
\[
\begin{aligned}
16\log\frac{f(\sqrt{-105})}{\sqrt[4]{2}} &= K(-3)\,K(140)\log\frac{T+U\sqrt{140}}{2}, \\
&\quad + K(5)\,K(-84)\log\frac{T+U\sqrt{84}}{2}, \\
&\quad + K(21)\,K(-20)\log\frac{T+U\sqrt{20}}{2}, \\
&\quad + K(-35)\,K(12)\log\frac{T+U\sqrt{12}}{2}.
\end{aligned}
\]
